\scanpage{57}
% archive foliation: 56
\begin{center}- 11 -\end{center}
\begin{flushright}n°385\end{flushright}
$\Psi_i$ et $\Psi'_i$ définies et continues dans un voisinage du point $\xi=(x_1(a),\ldots,x_n(a))$ de $K^n$ et telles que :
\[
\Phi=\sum_i\Psi_i\Psi'_i
\qquad \Psi_i(\xi)=\Psi'_i(\xi)=0
\]
et que les $\Psi_i$ soient de classe $C^r$. Ces conditions entraînent que
$\Phi(\lambda_1,\ldots,\lambda_n)/\sum_i|\lambda_i-\xi_i|$ tend vers $0$ si le point $(\lambda_1,\ldots,\lambda_n)$ tend vers $\xi$, et par conséquent que les dérivées de $\Phi$ sont nulles au point $\xi$. Comme on a par ailleurs $\langle\mu,f\rangle=0$ dans ce cas, la formule (24) est encore vraie.

La prop.1 montre que tout élément de $C_a^r(X)$ est congru modulo $I_a^2(X)$ au germe en $a$ d'un polynôme de degré $\leq1$ en les $x_i$ ; les cas particuliers précédemment démontrés prouvent que (24) est vraie en général.
\begin{flushright}C.Q.F.D.\end{flushright}

Considérons alors deux systèmes de coordonnées $x_1,\ldots,x_n$ et $y_1,\ldots,y_n$ au voisinage du point $a$ ; nous introduirons les fonctions de passage $\Phi_j$ telles que l'on ait :
\[
(26)\qquad y_j(p)=\Phi_j(x_1(p),\ldots,x_n(p))
\qquad 1\leq j\leq n
\]
lorsque $p$ appartient à l'intersection des domaines de définition des deux systèmes de coordonnées ; par ailleurs, on introduit la matrice jacobienne par :
\[
J_{ij}(p)=D_i\Phi_j(x_1(p),\ldots,x_n(p)).
\]
D'après le lemme 1 appliqué à $f=y_j$, on a alors :
\[
(27)\qquad
\langle\mu,y_j\rangle=\sum_{1\leq i\leq n}\langle\mu,x_i\rangle J_{ij}(a)
\qquad 1\leq j\leq n
\]
pour tout $\mu\in L$. D'après la définition d'un vecteur tangent, on a finalement la proposition suivante :

PROPOSITION 3.-- \uline{Soient} $a$ \uline{un point d'une variété} $X$ \uline{de classe} $C^r$ $(r\neq\unsure{\omega|0})$, \uline{et} $\bar D_a^1(X)$ \uline{l'espace des distributions} $\mu$ \uline{de support} $a$ \uline{et d'ordre} $\leq1$ \uline{telles que} $\langle\mu,1\rangle=0$.
\uline{Il existe un isomorphisme} $u$ \uline{de} $\bar D_a^1(X)$ \uline{sur l'espace tangent} $T_a(X)$ \uline{à} $X$ \uline{au point} $a$,
